\section{Loop 1: Build It Now}

\label{sec:loop1}

\subsection{Motivation and Trigger Condition}

The Build-It-Now (BIN) loop is the most novel contribution of this architecture.
Its design philosophy is that \emph{the best time to fix a recurring problem
is the third time it occurs}, not at the end of the session, not in the next
sprint, and not after it has been documented in a backlog. The third occurrence
provides enough evidence that the problem is systematic (ruling out one-off
errors) while the session context is still fresh enough to build a targeted
solution.

\begin{definition}[BIN Trigger]
\label{def:bin_trigger}
The BIN trigger $\mathcal{T}_{\text{BIN}}$ fires when the telemetry stream
contains a pattern match:
\begin{equation}
\mathcal{T}_{\text{BIN}}(\mathcal{T}_{\text{cur}}) = \mathbf{1}\bigl[
    \exists\, (t, f) \in \mathcal{V}_T \times \mathcal{V}_F :
    \#\{ \tau \in \mathcal{T}_{\text{cur}} \mid \tau.\text{tool}=t,\;
    \tau.\text{failure\_mode}=f \} \geq \theta
\bigr],
\end{equation}
where $\theta = 3$ is the default threshold and $\mathcal{T}_{\text{cur}}$
is the current session slice.
\end{definition}

The threshold $\theta = 3$ is motivated by the coupon-collector problem: with
$k$ distinct failure modes, the expected number of trials to observe any mode
three times is $3k$ when modes are uniformly distributed, and substantially
less when one mode is dominant (which is the case of interest). We show in
Section~\ref{sec:triggers} that $\theta \in \{2,3,4\}$ yields the best
precision-recall tradeoff, with $\theta = 3$ being Pareto-optimal in most
regimes.

\subsection{FRICTION-DETECT Algorithm}

Algorithm~\ref{alg:friction_detect} specifies the FRICTION-DETECT procedure
that continuously monitors the telemetry stream and fires the BIN trigger.

\begin{algorithm}[H]
\caption{FRICTION-DETECT$(S, \theta)$}
\label{alg:friction_detect}
\begin{algorithmic}[1]
\Require Telemetry stream $S$ (NATS subscription), threshold $\theta$
\Ensure BIN trigger events with $(t, f, \text{evidence})$ payload
\State $C \gets \textsc{DefaultDict}(\mathbb{N})$ \Comment{Failure counter: $(t,f) \to \mathbb{N}$}
\State $E \gets \textsc{DefaultDict}(\text{list})$ \Comment{Evidence collector: $(t,f) \to [\tau]$}
\State $\text{fired} \gets \emptyset$ \Comment{Set of already-triggered pairs}
\For{each record $\tau$ from $S$}
    \If{$\tau.\text{outcome} = \textsc{failure}$ \textbf{and} $\tau.\text{failure\_mode} \neq \bot$}
        \State $k \gets (\tau.\text{tool},\; \tau.\text{failure\_mode})$
        \State $C[k] \mathrel{+}= 1$
        \State $E[k].\textsc{append}(\tau)$
        \If{$C[k] \geq \theta$ \textbf{and} $k \notin \text{fired}$}
            \State $\text{fired}.\textsc{add}(k)$
            \State \textbf{emit} $\textsc{BINEvent}(t=k_1,\; f=k_2,\; \text{evidence}=E[k])$
        \EndIf
    \EndIf
\EndFor
\end{algorithmic}
\end{algorithm}

\subsection{BIN Build Pipeline}

When a BINEvent fires, the agent executes the following multi-stage pipeline.
The pipeline is designed to be atomic: if any stage fails, the agent rolls back
to its pre-BIN state and resumes the paused task.

\begin{enumerate}[leftmargin=1.5em]
    \item \textbf{Pause.} Serialize the current task state (context snapshot,
          pending actions queue) to a checkpoint file.

    \item \textbf{Design.} Prompt the LLM with: (a) the BINEvent evidence
          (up to 3 example failure records), (b) the current tool library
          manifest, and (c) a structured template requesting a micro-tool
          design. The design must specify: tool name, input schema, output
          schema, implementation sketch, and acceptance criteria.

    \item \textbf{Build.} Generate tool implementation from the design. Tools
          are TypeScript modules conforming to the Hanzo bot plugin interface.
          The LLM generates the implementation; a static linter and type
          checker run immediately.

    \item \textbf{Test.} Execute the acceptance criteria as unit tests against
          the reproduced failure case (reconstructed from the evidence records).
          A tool that does not pass its own acceptance criteria is not loaded.

    \item \textbf{Gate.} If all acceptance criteria pass, the tool is added to
          the tool library with a \texttt{status: "probationary"} flag. A
          probationary tool requires two more successful uses before its status
          is promoted to \texttt{"stable"}.

    \item \textbf{Reload.} Load the new tool into the active tool registry.

    \item \textbf{Resume.} Restore the serialized task state and continue
          execution with the new tool available.
\end{enumerate}

\begin{algorithm}[H]
\caption{BIN-BUILD-PIPELINE$(\text{event}, \mathcal{K})$}
\label{alg:bin_build}
\begin{algorithmic}[1]
\Require BINEvent $\text{event}$, current tool library $\mathcal{K}$
\Ensure Updated tool library $\mathcal{K}'$ or $\mathcal{K}$ on failure
\State $\text{checkpoint} \gets \textsc{SerializeContext}()$
\State $\text{design} \gets \textsc{LLM-Design}(\text{event.evidence}, \mathcal{K})$
\If{$\text{design} = \bot$}
    \State \Return $\mathcal{K}$ \Comment{LLM could not produce design; resume without change}
\EndIf
\State $\text{impl} \gets \textsc{LLM-Implement}(\text{design})$
\State $\text{lint\_ok} \gets \textsc{TypeCheck}(\text{impl})$
\If{$\neg \text{lint\_ok}$}
    \State \Return $\mathcal{K}$
\EndIf
\State $\text{test\_result} \gets \textsc{RunAcceptanceCriteria}(\text{impl}, \text{design.criteria}, \text{event.evidence})$
\If{$\neg \text{test\_result.passed}$}
    \State \Return $\mathcal{K}$
\EndIf
\State $\text{tool} \gets \textsc{Package}(\text{impl}, \text{status}=\texttt{"probationary"})$
\State $\mathcal{K}' \gets \mathcal{K} \cup \{\text{tool}\}$
\State $\textsc{ReloadRegistry}(\mathcal{K}')$
\State $\textsc{RestoreContext}(\text{checkpoint})$
\State \Return $\mathcal{K}'$
\end{algorithmic}
\end{algorithm}

\subsection{Anti-Deferral Properties}

\begin{proposition}[No Backlog Accumulation]
\label{prop:no_backlog}
In a system running the BIN protocol, the expected size of the improvement
backlog grows as $O(\log s)$ in the number of sessions $s$, compared to
$O(s)$ in deferral-based systems.
\end{proposition}

\begin{proof}
In a deferral-based system, each session generates at most $B$ improvement
opportunities (where $B$ is session length divided by minimum failure stride).
These are appended to the backlog. Without a consumption mechanism of comparable
rate, the backlog grows as $O(B \cdot s) = O(s)$.

In the BIN system, each failure pattern $(t, f)$ triggers a build the first
time it reaches threshold $\theta$. Subsequent occurrences of the same
$(t, f)$ pair after a successful build do not trigger another build (the tool
is already in the library). New $(t, f)$ pairs can only arise from the
vocabulary expansion rate, which grows at most as $O(\log s)$ by the
information-theoretic lower bound on new information in a stationary environment
(new failure modes are progressively rarer as the tool library matures).
Therefore the backlog size is $O(\log s)$.
\end{proof}

% ============================================================================
% 6. LOOP 2: ACTIVE LEARNING
% ============================================================================
